Imagine a launch with initial Sun-relative speed slightly lower than the
Earth's orbital speed. The spacecraft follows an ellipse with the Sun at
the focus opposite the launch point. As the initial speed decreases, the
ellipse becomes increasingly elongated. For an orbit tangent to the Sun,
\[
r_{\min}=R_S\approx\frac{r_{ES}}{200}.
\]
Its semiminor axis is negligible compared with the semimajor axis. The
trajectory can therefore be approximated by a straight free fall to the Sun.
The initial Sun-relative speed must be zero, so the Earth-relative launch
velocity is opposite to the Earth's orbital velocity with the same magnitude,
\[
v_{ES}=\frac{2\pi r_{ES}}{T_E}
=\frac{2\times3.14\times1.5\times10^{11}}{3.15\times10^7}
\approx3\times10^4\,\mathrm{m\,s^{-1}}=30\,\mathrm{km\,s^{-1}}.
\]
\ioaafig{2022_TH_Q08-s1.pdf}
Kepler's third law gives
\[
\frac{T_E^2}{r_{ES}^3}=\frac{T_{NC}^2}{(r_{ES}/2)^3},\qquad
T_{NC}=\frac{T_E}{2\sqrt2}.
\]
The fall is half the period,
\[t=\frac{T_{NC}}2=\frac{T_E}{4\sqrt2}\approx5.6\times10^6\,\mathrm{s}.\]
The answer is (c).
